% Practice Test 20 — Alaska Edition
% Questions: 30 (21 MC + 9 SA)
% Generated by scripts/generate_practice_tests.py

\begin{practiceQuestion}{1-3-q25}{sa}
\begin{questionText}
A factory produces widgets at a constant rate. In $2\frac{1}{2}$ hours it makes $175$ widgets. What is $k$ in widgets per hour?
\end{questionText}
\correctAnswer{$k = 70$ widgets per hour}
\explanation{$k = \frac{175}{2.5} = 70$ widgets per hour.}
\end{practiceQuestion}

\begin{practiceQuestion}{1-5-q02}{mc}
\begin{questionText}
A proportional graph passes through $(3, 18)$. What is the unit rate?
\end{questionText}
\choiceA{$3$}
\choiceB{$6$}
\choiceC{$15$}
\choiceD{$18$}
\correctAnswer{B}
\explanation{Unit rate $= k = \frac{18}{3} = 6$.}
\end{practiceQuestion}

\begin{practiceQuestion}{1-6-q17}{mc}
\begin{questionText}
A $12$-oz box of cereal costs $\$3.60$. A $20$-oz box costs $\$5.40$. Which is the better value?
\end{questionText}
\choiceA{$12$-oz box at $\$0.30$/oz}
\choiceB{$20$-oz box at $\$0.27$/oz}
\choiceC{$12$-oz box at $\$0.27$/oz}
\choiceD{They are the same price per ounce}
\correctAnswer{B}
\explanation{$12$-oz: $3.60 \div 12 = \$0.30$/oz. $20$-oz: $5.40 \div 20 = \$0.27$/oz. The $20$-oz box is the better value.}
\end{practiceQuestion}

\begin{practiceQuestion}{2-2-q28}{gmc}
\begin{questionText}
Look at the proportion model below. What value of $n$ makes the proportion true?

\begin{center}
\begin{tikzpicture}[scale=0.7]
  % Two bars representing the proportion
  % Top bar: n out of 80
  \draw[thick] (0,1.5) rectangle (8,2.3);
  \fill[funGreen!40] (0,1.5) rectangle (3,2.3);
  \draw[thick] (0,1.5) rectangle (8,2.3);
  \draw[thick] (3,1.5) -- (3,2.3);
  \node[above,font=\small] at (1.5,2.3) {$n$};
  \node[above,font=\small] at (5.5,2.3) {};
  \node[right,font=\small] at (8.2,1.9) {$80$};
  % Bottom bar: percent out of 100
  \draw[thick] (0,0) rectangle (10,0.8);
  \fill[funGreen!40] (0,0) rectangle (3.75,0.8);
  \draw[thick] (0,0) rectangle (10,0.8);
  \draw[thick] (3.75,0) -- (3.75,0.8);
  \node[below,font=\small] at (1.875,0) {$37.5$};
  \node[right,font=\small] at (10.2,0.4) {$100$};
  % Connection
  \node[font=\small] at (-0.5,1.15) {$=$};
\end{tikzpicture}
\end{center}
\end{questionText}
\choiceA{$28$}
\choiceB{$30$}
\choiceC{$32$}
\choiceD{$37.5$}
\correctAnswer{B}
\explanation{$\frac{n}{80} = \frac{37.5}{100}$. Cross-multiply: $100n = 3000$. Divide: $n = 30$.}
\end{practiceQuestion}

\begin{practiceQuestion}{2-5-q05}{mc}
\begin{questionText}
A real estate agent earns $3\%$ commission on a $\$250{,}000$ home sale. What does she earn?
\end{questionText}
\choiceA{$\$750$}
\choiceB{$\$2{,}500$}
\choiceC{$\$7{,}500$}
\choiceD{$\$75{,}000$}
\correctAnswer{C}
\explanation{$250{,}000 \times 0.03 = \$7{,}500$.}
\end{practiceQuestion}

\begin{practiceQuestion}{3-1-q23}{sa}
\begin{questionText}
Find two integers whose absolute values are both $7$.
\end{questionText}
\correctAnswer{$7$ and $-7$}
\explanation{$|7| = 7$ and $|-7| = 7$. A number and its opposite share the same absolute value.}
\end{practiceQuestion}

\begin{practiceQuestion}{3-2-q03}{mc}
\begin{questionText}
What is $15 + (-15)$?
\end{questionText}
\choiceA{$30$}
\choiceB{$-30$}
\choiceC{$15$}
\choiceD{$0$}
\correctAnswer{D}
\explanation{A number plus its opposite (additive inverse) always equals $0$.}
\end{practiceQuestion}

\begin{practiceQuestion}{3-3-q18}{mc}
\begin{questionText}
The high temperature was $4°$F and the low was $-11°$F. What is the difference between the high and low?
\end{questionText}
\choiceA{$7°$F}
\choiceB{$-7°$F}
\choiceC{$15°$F}
\choiceD{$-15°$F}
\correctAnswer{C}
\explanation{$4 - (-11) = 4 + 11 = 15°$F.}
\end{practiceQuestion}

\begin{practiceQuestion}{3-10-q15}{mc}
\begin{questionText}
A family's monthly budget is $\$3{,}200$. They spend $\dfrac{1}{4}$ on rent, $\dfrac{3}{16}$ on food, and $\$540$ on other expenses. How much is left?
\end{questionText}
\choiceA{$\$860$}
\choiceB{$\$1{,}060$}
\choiceC{$\$540$}
\choiceD{$\$1{,}260$}
\correctAnswer{A}
\explanation{Rent: $\frac{1}{4} \times 3{,}200 = \$800$. Food: $\frac{3}{16} \times 3{,}200 = \$600$. Total spent: $800 + 600 + 540 = \$1{,}940$. Left: $3{,}200 - 1{,}940 = \$1{,}260$.}
\end{practiceQuestion}

\begin{practiceQuestion}{4-1-q23}{sa}
\begin{questionText}
A cell phone plan charges $c$ cents per minute. Write an expression for the cost of a $15$-minute call in dollars.
\end{questionText}
\correctAnswer{$\frac{15c}{100}$ or $0.15c$}
\explanation{$15$ minutes at $c$ cents each is $15c$ cents. To convert to dollars, divide by $100$: $\frac{15c}{100} = 0.15c$.}
\end{practiceQuestion}

\begin{practiceQuestion}{5-2-q19}{sa}
\begin{questionText}
Solve $2x + 15 = 31$.
\end{questionText}
\correctAnswer{$x = 8$}
\explanation{Subtract $15$: $2x = 16$. Divide by $2$: $x = 8$. Check: $2(8) + 15 = 31$ \checkmark.}
\end{practiceQuestion}

\begin{practiceQuestion}{5-3-q04}{mc}
\begin{questionText}
Solve $-4(y + 3) = 28$.
\end{questionText}
\choiceA{$y = 4$}
\choiceB{$y = -10$}
\choiceC{$y = 10$}
\choiceD{$y = -4$}
\correctAnswer{B}
\explanation{Divide by $-4$: $y + 3 = -7$. Subtract $3$: $y = -10$. Check: $-4(-10 + 3) = -4(-7) = 28$ \checkmark.}
\end{practiceQuestion}

\begin{practiceQuestion}{5-4-q22}{sa}
\begin{questionText}
Solve $-0.3x + 4.2 = 1.8$.
\end{questionText}
\correctAnswer{$x = 8$}
\explanation{Subtract $4.2$: $-0.3x = -2.4$. Divide by $-0.3$: $x = 8$. Check: $-0.3(8) + 4.2 = -2.4 + 4.2 = 1.8$ \checkmark.}
\end{practiceQuestion}

\begin{practiceQuestion}{5-5-q11}{mc}
\begin{questionText}
A movie theater requires you to be more than $48$ inches tall for a ride. You are $42$ inches tall and grow $g$ inches per year. Which inequality finds how many years until you are tall enough?
\end{questionText}
\choiceA{$42 + g > 48$}
\choiceB{$42g > 48$}
\choiceC{$48 - g > 42$}
\choiceD{$42 - g > 48$}
\correctAnswer{A}
\explanation{Current height plus growth: $42 + g > 48$.}
\end{practiceQuestion}

\begin{practiceQuestion}{5-6-q07}{mc}
\begin{questionText}
The graph of an inequality has an open circle at $-1$ and is shaded to the left. Which inequality matches?
\end{questionText}
\choiceA{$x \leq -1$}
\choiceB{$x < -1$}
\choiceC{$x \geq -1$}
\choiceD{$x > -1$}
\correctAnswer{B}
\explanation{Open circle means $-1$ is not included, shading left means less than: $x < -1$.}
\end{practiceQuestion}

\begin{practiceQuestion}{6-2-q21}{sa}
\begin{questionText}
A blueprint uses scale $1:80$. You redraw at scale $1:20$. A room is $2$ cm by $3$ cm on the original. What is the area of the room on the new drawing?
\end{questionText}
\correctAnswer{$96$ cm$^2$}
\explanation{Scale ratio: $\frac{80}{20} = 4$. New dimensions: $2 \times 4 = 8$ cm and $3 \times 4 = 12$ cm. New area: $8 \times 12 = 96$ cm$^2$.}
\end{practiceQuestion}

\begin{practiceQuestion}{6-3-q16}{mc}
\begin{questionText}
A parallelogram has side lengths $5$ cm and $8$ cm. How many different parallelograms can be drawn?
\end{questionText}
\choiceA{None}
\choiceB{Exactly one}
\choiceC{Exactly two}
\choiceD{More than one}
\correctAnswer{D}
\explanation{Two side lengths of a parallelogram do not fix the angles. The shape can be a rectangle, a rhomboid, or anything in between. Many different parallelograms are possible.}
\end{practiceQuestion}

\begin{practiceQuestion}{6-4-q03}{mc}
\begin{questionText}
Can a triangle have angles $50°$, $60°$, and $80°$?
\end{questionText}
\choiceA{Yes, because the angles are all positive}
\choiceB{Yes, because each angle is less than $180°$}
\choiceC{No, because $50 + 60 + 80 = 190 \neq 180$}
\choiceD{No, because the angles are not equal}
\correctAnswer{C}
\explanation{The three angles sum to $190°$, not $180°$. Triangle angles must always sum to exactly $180°$.}
\end{practiceQuestion}

\begin{practiceQuestion}{6-5-q11}{mc}
\begin{questionText}
Which 3D figure always has circular cross-sections when cut parallel to its base?
\end{questionText}
\choiceA{Rectangular prism}
\choiceB{Triangular pyramid}
\choiceC{Cone}
\choiceD{Cube}
\correctAnswer{C}
\explanation{A cone has a circular base, and all horizontal cuts parallel to the base produce circles of varying sizes.}
\end{practiceQuestion}

\begin{practiceQuestion}{7-3-q14}{mc}
\begin{questionText}
A sprinkler waters a circular area with a radius of $5$ meters. How many square meters does it water? Use $\pi \approx 3.14$.
\end{questionText}
\choiceA{$31.4$ m$^2$}
\choiceB{$50$ m$^2$}
\choiceC{$78.5$ m$^2$}
\choiceD{$157$ m$^2$}
\correctAnswer{C}
\explanation{$A = \pi r^2 = 3.14 \times 25 = 78.5$ m$^2$.}
\end{practiceQuestion}

\begin{practiceQuestion}{7-4-q22}{sa}
\begin{questionText}
An H-shaped figure is made of three rectangles. The top and bottom bars are each $10$ cm by $2$ cm. The middle vertical bar is $2$ cm by $6$ cm. What is the total area?
\end{questionText}
\correctAnswer{$52$ cm$^2$}
\explanation{Top: $10 \times 2 = 20$ cm$^2$. Bottom: $10 \times 2 = 20$ cm$^2$. Middle: $2 \times 6 = 12$ cm$^2$. Total: $20 + 20 + 12 = 52$ cm$^2$.}
\end{practiceQuestion}

\begin{practiceQuestion}{7-5-q04}{mc}
\begin{questionText}
How many faces does a rectangular prism have?
\end{questionText}
\choiceA{$4$}
\choiceB{$5$}
\choiceC{$6$}
\choiceD{$8$}
\correctAnswer{C}
\explanation{A rectangular prism has $6$ faces: top, bottom, front, back, left, and right.}
\end{practiceQuestion}

\begin{practiceQuestion}{7-6-q17}{mc}
\begin{questionText}
A moving box is $2$ ft long, $2$ ft wide, and $3$ ft tall. How many of these boxes can fit in a space that is $6$ ft by $4$ ft by $3$ ft?
\end{questionText}
\choiceA{$3$}
\choiceB{$6$}
\choiceC{$12$}
\choiceD{$18$}
\correctAnswer{B}
\explanation{Space volume: $6 \times 4 \times 3 = 72$ ft$^3$. Box volume: $2 \times 2 \times 3 = 12$ ft$^3$. Number: $72 \div 12 = 6$ boxes.}
\end{practiceQuestion}

\begin{practiceQuestion}{8-1-q01}{mc}
\begin{questionText}
A school has $800$ students. The principal surveys $50$ students about their favorite subject. What is the population?
\end{questionText}
\choiceA{The $50$ surveyed students}
\choiceB{All $800$ students in the school}
\choiceC{The principal}
\choiceD{The students who like math}
\correctAnswer{B}
\explanation{The population is the entire group you want to study — all $800$ students.}
\end{practiceQuestion}

\begin{practiceQuestion}{8-3-q06}{mc}
\begin{questionText}
Team A's scores: mean $72$, MAD $3$. Team B's scores: mean $74$, MAD $8$. Which team is more consistent?
\end{questionText}
\choiceA{Team A because its MAD is smaller}
\choiceB{Team B because its mean is higher}
\choiceC{Team A because its mean is lower}
\choiceD{Team B because its MAD is larger}
\correctAnswer{A}
\explanation{A smaller MAD means scores are closer to the mean, so Team A is more consistent.}
\end{practiceQuestion}

\begin{practiceQuestion}{8-4-q12}{mc}
\begin{questionText}
Two groups have the same mean but different MADs. What does this tell you?
\end{questionText}
\choiceA{One group scored higher than the other}
\choiceB{The groups have different levels of variability}
\choiceC{The groups are identical}
\choiceD{The means were calculated incorrectly}
\correctAnswer{B}
\explanation{Same mean means same center. Different MADs mean one group's data is more spread out than the other's.}
\end{practiceQuestion}

\begin{practiceQuestion}{8-7-q20}{sa}
\begin{questionText}
A line graph shows monthly rainfall: Jan $2$ in., Feb $3$ in., Mar $5$ in., Apr $4$ in. Find the total rainfall and the month with the most.
\end{questionText}
\correctAnswer{Total $= 14$ in.; March had the most ($5$ in.)}
\explanation{$2+3+5+4 = 14$ inches. March has the highest value at $5$ in.}
\end{practiceQuestion}

\begin{practiceQuestion}{9-4-q17}{mc}
\begin{questionText}
A bag has tiles labeled A, A, A, B, B. What is the probability of drawing a B?
\end{questionText}
\choiceA{$\frac{1}{5}$}
\choiceB{$\frac{2}{3}$}
\choiceC{$\frac{2}{5}$}
\choiceD{$\frac{3}{5}$}
\correctAnswer{C}
\explanation{There are $2$ B tiles out of $5$ total. $P(B) = \frac{2}{5}$.}
\end{practiceQuestion}

\begin{practiceQuestion}{9-6-q29}{gsa}
\begin{questionText}
The table shows all possible products when two dice are rolled. How many outcomes give a product of $6$? Write the probability as a fraction in simplest form.

\begin{center}
\begin{tabular}{|c|c|c|c|c|c|c|}
\hline
$\times$ & \textbf{1} & \textbf{2} & \textbf{3} & \textbf{4} & \textbf{5} & \textbf{6} \\
\hline
\textbf{1} & $1$ & $2$ & $3$ & $4$ & $5$ & $6$ \\
\hline
\textbf{2} & $2$ & $4$ & $6$ & $8$ & $10$ & $12$ \\
\hline
\textbf{3} & $3$ & $6$ & $9$ & $12$ & $15$ & $18$ \\
\hline
\textbf{4} & $4$ & $8$ & $12$ & $16$ & $20$ & $24$ \\
\hline
\textbf{5} & $5$ & $10$ & $15$ & $20$ & $25$ & $30$ \\
\hline
\textbf{6} & $6$ & $12$ & $18$ & $24$ & $30$ & $36$ \\
\hline
\end{tabular}
\end{center}
\end{questionText}
\correctAnswer{$\frac{1}{9}$}
\explanation{Products of $6$: $(1,6), (2,3), (3,2), (6,1)$ — $4$ outcomes. $P = \frac{4}{36} = \frac{1}{9}$.}
\end{practiceQuestion}

\begin{practiceQuestion}{9-7-q16}{mc}
\begin{questionText}
Which of the following is the FIRST step in designing a simulation?
\end{questionText}
\choiceA{Record the results}
\choiceB{Run many trials}
\choiceC{Identify the event and its possible outcomes}
\choiceD{Calculate the experimental probability}
\correctAnswer{C}
\explanation{Before choosing a model or running trials, you must first identify what event you're simulating and what the possible outcomes are.}
\end{practiceQuestion}
